\section{Anexo C: Derechos de autor y propiedad intelectual}

El desarrollo de una plataforma integral como SocioUnido, que involucra arquitecturas distribuidas, algoritmos de predicción mediante inteligencia artificial y protocolos criptográficos de validación fuera de línea (\textit{offline}), constituye un activo tecnológico de alto valor. Por consiguiente, la protección de la propiedad intelectual y la definición clara de los derechos de explotación comercial resultan imperativas para resguardar el esfuerzo del equipo de ingeniería, sin obstaculizar el propósito educativo y colaborativo del Trabajo Profesional.

\subsection{Modelo de licenciamiento: PolyForm Noncommercial License 1.0.0}

Para el presente proyecto, se ha adoptado la licencia \textbf{PolyForm Noncommercial License 1.0.0}. La elección de este modelo legal responde a una necesidad dual: por un lado, blindar el software contra la explotación o monetización no autorizada por parte de terceros comerciales; y por el otro, mantener el código fuente abierto y accesible para fomentar la transparencia, la revisión entre pares y el aprendizaje dentro del ecosistema universitario.

A diferencia de las licencias de código abierto (\textit{Open Source}) tradicionales (como MIT o Apache 2.0), que permiten el uso comercial irrestricto, la licencia PolyForm establece una barrera legal clara que prohíbe cualquier aplicación lucrativa sin el consentimiento previo, expreso y por escrito de los autores originales. 

\subsection{Condiciones y permisos estipulados}

El marco legal de PolyForm Noncommercial License 1.0.0 otorga a los usuarios que accedan al código fuente de SocioUnido una serie de libertades estrictamente delimitadas a fines no lucrativos:

\begin{itemize}
    \item \textbf{Uso personal y académico:} Se permite explícitamente la copia, instalación y ejecución del software para investigación, experimentación, pruebas y estudio personal. Esto resulta fundamental para que docentes, evaluadores y futuros estudiantes puedan auditar la plataforma, analizar su arquitectura y comprender sus flujos de integración de manera irrestricta.
    \item \textbf{Modificaciones y obras derivadas:} La licencia autoriza la realización de cambios en el código fuente y la creación de trabajos derivados, siempre y cuando estos esfuerzos se mantengan bajo el paraguas de propósitos no comerciales (por ejemplo, bifurcaciones o \textit{forks} con fines netamente educativos o de prueba).
    \item \textbf{Distribución restringida:} Se concede el derecho a distribuir copias del software o de sus trabajos derivados, condicionado a que los receptores de dichas copias queden obligatoriamente sujetos a los mismos términos legales y restricciones de uso comercial.
    \item \textbf{Organizaciones exentas:} El uso por parte de instituciones educativas, organizaciones sin fines de lucro o entidades de investigación pública es considerado un propósito permitido. Esta cláusula se alinea de manera perfecta con la presentación del sistema ante la Universidad de Buenos Aires (UBA), garantizando que la institución posea los permisos necesarios para interactuar con la plataforma.
\end{itemize}

\newpage

\subsection{Aviso legal requerido (\textit{Required notice})}

Para garantizar el cumplimiento de las normativas de propiedad intelectual y hacer efectiva la protección de los derechos de autor en cualquier distribución o inspección de los repositorios, se ha incrustado el siguiente aviso legal vinculante en el código fuente de SocioUnido. Todo tercero que obtenga una copia del software o sus repositorios está legalmente obligado a conservar intacta esta cláusula:

\begin{quote}
    \texttt{Required Notice: Copyright 2026 Ascencio Felipe Santino, Ghosn Lautaro Gabriel, Guerrero Martín, Zielonka Axel (Socio Unido). \\
    Any use of this software for commercial purposes is strictly prohibited without prior written consent from the copyright holders.}
\end{quote}

\subsection{Impacto y contribución al ámbito universitario}

La adopción de esta postura de licenciamiento refleja un compromiso profundo con los valores académicos de la Facultad de Ingeniería. Al permitir que el código fuente, los repositorios en GitHub, las métricas de calidad y los flujos de integración sean plenamente visibles y editables en un entorno de pruebas, SocioUnido se convierte en un caso de estudio real, táctico y auditable. 

La licencia \textit{PolyForm Noncommercial} garantiza que este valioso aporte al conocimiento público y a la comunidad universitaria permanezca intacto y protegido de la apropiación indebida. De esta manera, el equipo desarrollador retiene de manera exclusiva el control absoluto sobre los derechos de comercialización de la plataforma como software como servicio (B2B SaaS).
